% ==============================================================================
% 第 24 章：小结 (Chapter 24: Summary)
% ==============================================================================

\chapter{小结}
\label{sec:ch24_summary}

本书第四部分的核心主题是：\textbf{设计处理需要多个相互交织定义来描述的数据形态的函数}。在真实世界中，此类交织数据无处不在——从你计算机操作系统中的本地文件目录树，到万维网上的网页超文本与 XML/JSON 报文，再到动画电影制作中精密的几何层次模型。

在细致研读并实践了本部分的全部内容之后，你已经充分见证了程序设计配方在应对这些高度复杂的数据形态时所展现出的卓越适应性与强大威力：

\begin{enumerate}
  \item \textbf{相互递归与并行模板：} 当程序数据的描述需要多个相互引用的数据定义（互递归定义）时，设计配方要求我们\textbf{同步并行开发一组模板}，每一个模板专职对应一份数据定义。如果数据定义 A 在某个位置引用了数据定义 B，那么在函数 A 的模板中，完全相同的逻辑位置就会以完全相同的方式调用函数 B。在其余方面，设计配方依然一如既往地逐个函数稳步推进。数据模型与代码实现之间高度对称的内在结构，正是计算思维的优雅所在。

  \item \textbf{处理多复杂输入的分类法：} 当一个函数需要同时接收两个复杂类型的输入参数时，我们需要清晰界定三种情形：
  \begin{itemize}
    \item \textbf{情形 1：被动传递。} 函数仅对其中一个主控参数实施递归拆解，而将另一个复杂参数视为原子值不加破坏地逐层传递；
    \item \textbf{情形 2：齐头并进。} 两个输入参数在假定下拥有完全同步的平行结构，函数以步调一致的锁步方式同步拆解两个输入；
    \item \textbf{情形 3：全组合矩阵。} 两个输入参数互相独立且各自拥有多个结构分支。此时必须构建二维分析矩阵表格，枚举出所有可能的分支组合，再据此编写条件分支并化简冗余。
  \end{itemize}
\end{enumerate}

在极少数情况下，如果一个函数必须处理三个相互独立的复杂参数，你完全可以类推：在脑海中建立一个三维的分析立方体，逐格推导其处理逻辑。

到目前为止，你已经掌握了职业程序员在其整个职业生涯中所能遇到的\textbf{所有基本的结构化数据形态}。今后无论面对多么扑朔迷离的现实业务需求，只要你牢记程序设计配方的步骤，它就能为你拨云见日，指引你迈出坚实可靠的第一步。
